% STATUS: PLANNING MAP ONLY; all theorem nodes are TODO.
\section{Theorem dependency map and completion gates}
\label{sec:dependency-map}
Arrows mean a prerequisite proof module, not a theorem already proved.
Every module also depends on the canonical contract and notation. The
following displayed adjacency graph is the full direct-dependency graph;
transitive prerequisites are inherited. Conditional edges are listed separately.
\[
\begin{array}{rcl}
\text{Contract}&\longrightarrow&\mathrm{P01}\\
\mathrm{P01}&\longrightarrow&\mathrm{P02}\\
\mathrm{P01},\mathrm{P02}&\longrightarrow&\mathrm{P03}\\
\mathrm{P03}&\longrightarrow&\mathrm{P04}\\
\mathrm{P04}&\longrightarrow&\mathrm{P05}\\
\mathrm{P05}&\longrightarrow&\mathrm{P06}\\
\mathrm{P05}&\longrightarrow&\mathrm{P07}\\
\mathrm{P01},\mathrm{P02},\mathrm{P04},\mathrm{P05}&\longrightarrow&\mathrm{P08}\\
\mathrm{P02},\mathrm{P03},\mathrm{P04}&\longrightarrow&\mathrm{P09}\\
\mathrm{P02},\mathrm{P03},\mathrm{P04},\mathrm{P09}&\longrightarrow&\mathrm{P10}\\
\mathrm{P02},\mathrm{P04},\mathrm{P05}&\longrightarrow&\mathrm{P11}\\
\mathrm{P11}&\longrightarrow&\mathrm{P12}\\
\mathrm{P11},\mathrm{P12}&\longrightarrow&\mathrm{P13}\\
\mathrm{P11},\mathrm{P12},\mathrm{P13}&\longrightarrow&\mathrm{P14}\\
\mathrm{P02},\mathrm{P04}&\longrightarrow&\mathrm{P15}\\
\mathrm{P04},\mathrm{P15}&\longrightarrow&\mathrm{P16}\\
\mathrm{P04},\mathrm{P16}&\longrightarrow&\mathrm{P17}\\
\mathrm{P04},\mathrm{P12},\mathrm{P13},\mathrm{P15},\mathrm{P16}&\longrightarrow&\mathrm{P18}\\
\end{array}
\]
The graph is acyclic. P07 is an additional prerequisite when a five-node
specialization is used in P11--P14. An optional periodic-orbit persistence
extension in P15 would additionally require P14; it is not part of the core
equilibrium-persistence gate. Source verification is a conditional input to
any six-term extension in P08. An unresolved optional extension does not block
the core module if its statement is explicitly excluded from the completed scope.

\subsection{One dependency-respecting completion order}
\begin{enumerate}
\item P01: Well-posedness and attracting bounded region.
\item P02: Synchrony and complete synchronous-equilibrium classification.
\item P03: Origin and homogeneous linearization.
\item P04: Arbitrary-(n,q) Fourier spectrum and normality.
\item P05: Exact step-two mode selection.
\item P06: Large-n selected-wave-number limit.
\item P07: Five-node golden-ratio specialization.
\item P08: Matched bilinear analysis and exact comparison.
\item P09: Transcritical bifurcation at the origin.
\item P10: Saddle-node bifurcation and true bistable region.
\item P11: Hopf spectral conditions and spectral simplicity.
\item P12: Hopf transversality.
\item P13: First Lyapunov coefficient and nondegenerate Hopf theorem.
\item P14: Equivariant rotating waves.
\item P15: Persistence under parameter heterogeneity.
\item P16: Static sensitivity and localized Green function.
\item P17: Transient and harmonic response.
\item P18: Stability radius and Hopf-threshold sensitivity.
\end{enumerate}
After P05, P06, P07 and the later independent branches need not be worked in
strict numerical order. This list is one valid order for the compendium and
for eventual completion. P11 collects spectral conditions, P12 establishes
transversality, and P13 is the home of the assembled nondegenerate Hopf theorem.
P14 owns the additional equivariant rotating-wave conclusion.

\subsection{Numerical dependency gates}
\begin{longtable}{>{\raggedright\arraybackslash}>{\raggedright\arraybackslash}p{0.31\textwidth}>{\raggedright\arraybackslash}p{0.59\textwidth}}
\toprule Task & Required analytical inputs before the claim is validated \\
\midrule\endhead
N01 Jacobian and boundary checks & P01--P03; P08 for MB. \\
N01 equilibrium and spectrum checks & P02--P04; P05--P07 for selection, limits and specialization. \\
N01/N02 bifurcation validation & P09--P13; complete-spectrum, transversality and independent l1 checks. \\
N02 periodic branch and wave validation & P13--P14; use Floquet and phase checks, not a time-series image alone. \\
N01/N02 heterogeneity validation & P15--P16 for equilibrium/remainder; P18 for threshold shifts. \\
N01/N03 response and robustness & P16--P18; specify norms, forcing and direct versus equilibrium-induced perturbations. \\
N03 final figures & Each figure's listed proof owners, passed N01 checks and reproducible N02 output where relevant. \\
\bottomrule
\end{longtable}
Exploratory future checks may help find errors in DRAFT derivations, but cannot
promote a theorem or a numerical claim to verified evidence on their own.
No computational checks are performed in this initialization.

\subsection{Manuscript claims withheld until verification}
\begin{longtable}{>{\raggedright\arraybackslash}>{\raggedright\arraybackslash}p{0.41\textwidth}>{\raggedright\arraybackslash}p{0.49\textwidth}}
\toprule Withheld claim & Minimum proof gate \\
\midrule\endhead
Local well-posedness, positivity, global existence and bounded attraction & P01 with precise attraction/absorption language. \\
Complete synchronous-equilibrium classification and origin stability & P02--P03; nonlinear boundary cases handled separately. \\
Arbitrary-(n,q) spectrum and normality & P04. \\
Exact q=2 selection, large-n limit and five-node constants & P05, P06 and P07, respectively. \\
Exact matched-bilinear comparison & P08 plus its prerequisites and an explicit calibration for quantitative comparisons. \\
Transcritical, fold and true bistability claims & P09--P10 with all-mode stability and explicit basin scope. \\
Nondegenerate symmetry-breaking Hopf & P11--P13, including rigorous nonzero frequency, simplicity, transversality and l1 nonvanishing. \\
Stable rotating-wave branch & P14 and applicable transverse stability from P11--P13. \\
Persistence, sensitivity and localized response & P15--P16; positivity, invertibility and stability assumptions distinguished. \\
Transient/harmonic response and robustness & P17--P18, with explicit norms and remainder bounds. \\
Numerical values, periods, figures or observed regimes & Applicable verified proofs plus N01/N02 evidence and recorded reproduction. \\
Novelty, priority or literature comparisons & Completed source-specific literature audits; none confirmed here. \\
\bottomrule
\end{longtable}
Before submission, mathematical claims must also have a completed independent
audit for the exact version cited. Unsupported candidate claims remain in the
archive and registry. No final-paper narrative is supplied by this map.

% Phase 1C planning overlay; no theorem or proof status is promoted.
\subsection{Phase 1C frozen scope and current attempt order}
The adjacency graph above remains in force. The numerical ordering above and
in the compendium is a valid reading order. The current research attempt order,
superseding earlier scheduling recommendations, is
\[
\begin{gathered}
\mathrm{P01}\to\mathrm{P02}\to\mathrm{P03}\to\mathrm{P04}\to\mathrm{P05}
\to\mathrm{P08}\to\mathrm{P07}\to\mathrm{P09}\to\mathrm{P10},\\
\mathrm{P11}\to\mathrm{P12}\to\mathrm{P13}\to\mathrm{P14}\to\mathrm{P06},\\
\mathrm{P15}\to\mathrm{P16}\to\mathrm{P17}\to\mathrm{P18}
\quad\text{(postponed extension queue).}
\end{gathered}
\]
Arrows in this schedule indicate order of attempts, not additional mathematical
dependencies. P06 is attempted only if it adds a useful large-$n$ consequence.
P07 precedes any five-node use in P11--P14. P15--P18 are outside the frozen paper.

\begin{longtable}{p{0.32\textwidth}p{0.24\textwidth}p{0.34\textwidth}}
\toprule Intended claim & Assembly gate & Frozen role \\
\midrule\endhead
Realizable $q=2$ selection and switching & P05 & Main theorem target; quadratic maximization alone is insufficient. \\
Matched DM--MB contrast & P08; P13 if comparing criticality & Main theorem target on a common admissible branch domain. \\
First symmetry-breaking Hopf & P11, P12, P13 & Assembled main theorem; P11 alone is spectral support. \\
Nontrivial $l_1$ sign region & P13 & Main theorem target; currently unresolved. \\
Coefficient reduction & P13 & Proposition/lemma; detailed algebra may be supplementary. \\
Exact rotating-wave branch & P14 & Corollary of verified Hopf and symmetry hypotheses. \\
Five-node motivating application & P07 plus applicable P11--P14 & Corollary within the main paper. \\
General spectrum; branch classification and local branch bifurcations & P01--P04, P09--P10 & Supporting propositions/lemmas. \\
Useful large-$n$ consequence & P06 & Conditional corollary. \\
Heterogeneity, response and robustness & P15--P18 & Postpone. \\
\bottomrule\end{longtable}
The primary comparison uses identical positive coefficient tuples with each
model's own named positive synchronous branch; branch/path choices must be
stated before their use. A criticality comparison is assembled after P13 and
does not introduce a circular prerequisite for P08's spectral analysis.

The numerical gates above remain valid. Checks and figures for P15--P18 are
dormant in this paper. No numerical task is executed in Phase 1C. All proof
statuses remain TODO. Exact-model priority is not closed because the named
2021--2024 main-text equation/theorem inventories remain incomplete; this is
a separate publication-claim gate, not a theorem dependency. See
\texttt{frozen\_theorem\_package.md} and \texttt{final\_novelty\_verdict.md}.
